\subsection{An integral reduction from a pair of idempotents}

\begin{theorem}[Proved]
\label{CP2:A-idempotent-pair-integral-reduction}
Let $M$ be a finite monoid with distinct idempotents $a,e$ satisfying
$ae=ea=e$.  Suppose the equivalence relation identifying $a$ with $e$
and leaving every other element in its own class is a monoid
congruence.  Put $Q=M/(a=e)$, and suppose that $h=a-e$ is central
in $\mathbb ZM$.

For every finite nonempty left-$N$, right-$M$ biset $S$, there is
an actual functor from its two-object collage $C$ to a collage $D$
with corners $Q,N$ and at most $|S|$ connecting arrows, inducing
an integral homology equivalence and a surjection on fundamental
groups.  If $BC$ is simply connected, $BD\simeq BC$.
Consequently an eighteen-arrow projective-plane collage cannot have
such a source corner.  The opposite assertion applies at its sink.
\end{theorem}
\begin{proof}
The condition $e=1$ would force $a=e$ and is impossible.  We do
allow $a=1$; in that case the quotient identity is the class of
$1,e$.  The construction always uses the actual quotient identity,
so no choice of a unital submonoid section is required.

Let $\pi:M\to Q$ be the quotient and $R=\mathbb ZM$.
Its linear kernel is exactly $\mathbb Zh$, because only one pair
of basis elements is identified.  Direct calculation gives
\[
 h^2=a-ae-ea+e=a-e=h.
\]
Thus $h$ is a central idempotent and $p=1-h$ is its complement.
Since the kernel is already a two-sided ideal containing $h$,
$Rh=\mathbb Zh$.  Hence
\[
 R=Rp\times\mathbb Zh,
 \qquad Rp\xrightarrow[\pi]{\cong}\mathbb ZQ.
\]
The displayed isomorphism is integral and unital, with $p$ mapping
to the identity of $\mathbb ZQ$.  Augmentation kills $h$, so the
trivial left $R$ module belongs to the $p$ factor.

Put $A(s)=sa$ and $E(s)=se$.  The idempotents $A,E$ obey
$AE=EA=E$, and commute with every left $N$ actor.  Form $T$ by
identifying each $t\in AS$ with $Et$; its canonical representatives
are $(S\setminus AS)\amalg ES$.  This is $N$ equivariant.  The
linear quotient $q_S:\mathbb Z[S]\to\mathbb Z[T]$ has kernel
\[
 \operatorname{span}_{\mathbb Z}\{[As]-[Es]:s\in S\}
   =\mathbb Z[S]h.
\]
Centrality of $h$ makes this kernel stable under the right $M$
action.  Since distinct basis elements of $\mathbb Z[T]$ are
independent, stability means precisely that every right actor
preserves the state equivalence relation.  This gives an actual
right $M$ action on $T$.  Its actors $a,e$ agree by construction,
so it factors through the actual quotient monoid $Q$.  Together
with $\pi$, this produces a functor $q:C\to D$ which is the
identity on objects and is surjective on all hom-sets.

The restriction
\[
 \mathbb Z[S]p\longrightarrow\mathbb Z[T]
\]
is an integral bimodule isomorphism.  Explicitly its basis is
\[
 [s]-[As]+[Es]\quad(s\notin AS),\qquad [c]\quad(c\in ES),
\]
mapping to the representative-state basis of $T$.  For $t\in AS$
one has $[t]p=[Et]$, proving spanning.  The separate coefficients
on each state outside $AS$, followed by the coefficients on $ES$,
prove independence over $\mathbb Z$.

The central direct-product decomposition now gives the natural
integral comparison
\[
 \mathbb Z[S]\otimes_{\mathbb ZM}^{\mathbf L}\mathbb Z
 \simeq
 \mathbb Z[S]p\otimes_{Rp}^{\mathbf L}\mathbb Z
 \cong
 \mathbb Z[T]\otimes_{\mathbb ZQ}^{\mathbf L}\mathbb Z.
\]
As in Theorem~\ref{CP2:A-six-constant-idempotent-integral-reduction},
this is the actual quotient comparison and is natural over
$\mathbb ZN$.  Applying derived $N$ coinvariants proves that the
bridge map is an integral homology equivalence.  The same direct
factor proves this for $BM\to BQ$, while $BN$ is unchanged.
The actual two-block mapping-cylinder sequence therefore makes
$Bq$ an integral homology equivalence in every degree.

Every edge-path loop of $BD$ lifts to one of $BC$, because $q$ fixes
objects and is onto their hom-sets.  Thus its fundamental-group map
is onto.  If $BC$ is simply connected, so is $BD$, and integral
homological Whitehead proves the asserted homotopy equivalence.
Since $|Q|=|M|-1$ and $|T|\leq|S|$, it strictly decreases the arrow
count.  The eighteen-arrow consequence follows from
Theorem~\ref{CP2:A-global-lower-eighteen}.  Opposition gives the
sink version.
\end{proof}

\begin{remark*}
The table predicate tests every ordered pair of distinct idempotents:
$ae=ea=e$, stability of the sole two-element congruence class under
all left and right translations, and the literal integral identities
$(a-e)m=m(a-e)$ for every basis element $m$.  No rational linear
solver is required.  The constant-idempotent six-element source is
the special case with $e=e_{00}$.  The quotient need not embed as
a unital submonoid of $M$.
\end{remark*}
